% §5 The corners --- blueprint chapter 10 (blueprint/src/parts/10-corners.tex), statements and
% proofs of record verbatim. Drafted 2026-09-15, from the blueprint at the state after the step-2
% merges (R160: prop:student-t(3) conditional on the Thorin representation, lem:thorin-ggc
% narrowed, lem:thorin-two-sided new, rem:student-ggc new) and the two drafting-day repairs
% (R164: the tail sentence of prop:matern-density restored to the statement; R165: a proof
% environment for lem:thorin-two-sided).
%
% Prose rewritten 2026-09-18 at the author's review (statements and proofs untouched). The
% author's findings: four lines of introduction for the section that is the payoff of the
% theory; no map; every subsection opening with technicalities about a theorem and its proof
% the reader has not seen, and with what is proved in Lean in place of explanation; the corner
% families not placed in their literature; no illustrations. What was done: the section opens
% with what the three families are, where they are used, the five results and why they matter,
% and a figure of the three families by parameter (figures/fig-corners.png from
% scripts/make-fig-corners.py); each subsection opens with motivation or intuition and the
% literature of its family (Whittle, Guttorp-Gneiting, Porcu et al., Lindgren et al.,
% Madan-Seneta, Burgeth et al. for Matern; Thorin and Bondesson for the Thorin class; Grosswald,
% Ismail-Kelker, Halgreen, Barndorff-Nielsen-Halgreen, Felsberg-Sommer, Duits et al. for
% Student-t; Pauwels et al., Duits et al., Pedersen et al. for stable); comments on a statement
% or a proof come AFTER the statement, before the proof; the machine-checking talk is out of the
% body; rem:matern-literature is moved from the end of the section to the end of 5.2 and given
% its citekeys; rem:student-ggc's second paragraph is reduced to what a reader needs (both
% paper-side changes to copied remarks). MIRRORED 2026-09-19 (R174): rem:matern-literature's
% citekeys are in the blueprint; rem:student-ggc's reduction is NOT mirrored, deliberately, the
% blueprint keeping its account of what is not machine-checked and why A18 is unadmitted (D5).
%
% Contents and departures from the blueprint chapter:
%   - all eleven statement nodes transcribed under their markers: prop:moments-tails ([A] on A13
%     and A14, notready on the bounded-support and heavy-tail clauses), prop:matern-exponent
%     (also Appendix A.3 of the line paper, in the form that paper's body used; stated here in
%     full with the blueprint's own remark at the end of its proof), prop:matern-density ([A] on
%     A15 and A16, notready), thm:matern, cor:matern-realization (no Lean scheduled),
%     prop:thorin-subclass ([A] on A11), lem:thorin-two-sided, lem:thorin-ggc ([A] on A17,
%     unadmitted, notready), prop:student-t (notready on clause (4) alone), prop:stable-family;
%     the [A] nodes' page anchors are printed in the prose before each statement, read from
%     blueprint/AXIOMS.md;
%   - rem:student-ggc and rem:matern-literature are copied (remarks carry no marker);
%     rem:student-ggc is where the one cited fact about the Student-t corner binds;
%   - the paragraph after prop:student-t on which locality principle cuts the Student-t family
%     out (module C's chapters) becomes a sentence naming the next module's subject with no
%     statement number;
%   - the blueprint's status annotations are not transcribed; what a reader needs from them is
%     said in prose next to each statement, in particular which clauses are machine-checked and
%     on what, and what a statement does not claim;
%   - this section's one relation-to-the-causal-theory remark (ADR-0003) is the paragraph after
%     the Thorin proposition on the bijection of the corners; the causal article is named there.

\section{The corners: \texorpdfstring{\Matern{}}{Matern}, Student-t, symmetric stable}
\label{sec:corners}

The two preceding sections described the admissible cone as a whole: its extreme rays, and the
part of it that the causal theory reaches. This section is about individual members, the ones
a user of scale space would compute with. Three families are treated, each already in use.
The \emph{\Matern{}} kernels are the
covariance family of spatial statistics (\sscite{whittle1954stationary};
\sscite{guttorp2006studies} for the history; \sscite{porcu2024matern} for the present state).
The \emph{Student-t} laws are the classical heavy-tailed relatives of the normal law, and
their member with one degree of freedom is the Cauchy kernel of the Poisson scale space
(\sscite{felsberg2004monogenic}). The \emph{symmetric stable} laws are the kernels the
semigroup axiomatics of scale space arrive at (\sscite{pauwels1995extended};
\sscite{duits2004axioms}). What the hemigroup theory adds is one description for all three and
a way to compare them. Each family is a curve of profiles: an exponential for the \Matern{}
kernels, the Gaussian mixture of the causal Bessel catalogue for the Student-t laws, a pure
power for the stable laws. The profile answers the practical questions about a kernel at
once: how many moments it has and how heavy its tail is, whether it is
bounded at the origin (Corollary~\ref{cor:origin-boundedness}), and whether it can be computed
exactly by recursive filters (\S\ref{sec:implementation}). Figure~\ref{fig:corners} shows the
three families.

The section has five parts. \S\ref{sec:moments-tails} gives the two dials on which the families
are compared: the $n$-th moment is finite exactly when $\int_1^\infty x^{n-1}k(x)\,dx$ is, and
the tail is governed by the support of the profile (Proposition~\ref{prop:moments-tails}).
\S\ref{sec:matern} is the main result: four descriptions single out the \Matern{} family, and
its members of integer index are exactly the families computed by a cascade of identical
first-order recursive sections (Theorem~\ref{thm:matern},
Corollary~\ref{cor:matern-realization}). \S\ref{sec:thorin} treats the profiles that are
mixtures of exponentials, the subcone that contains all three families and on which the
bridge of \S\ref{sec:bridge} is a bijection (Proposition~\ref{prop:thorin-subclass}).
\S\ref{sec:student} and \S\ref{sec:stable} place the Student-t and the stable families in it
(Propositions~\ref{prop:student-t} and~\ref{prop:stable-family}). Table~\ref{tab:families}
collects what the section and the two around it establish for each family, with the unit-step
members of \S\ref{sec:cone} for contrast.

\begin{table}[!htb]
\centering
\caption{The named families compared, at canonical scale $t$ and with no Gaussian part
except in the first row. ``Lag sections'' means an exact realization of every stage by the
first-order forward--backward sections of \S\ref{sec:implementation}; ``not a limit'' means
that no sequence of families so realized converges to the member, by
Proposition~\ref{prop:rational-closure} --- at its clause (3) for the unit-step members and the
\Matern{} members of non-integer index, and at its clause (5), on the moments alone, for the
stable and Student-t members. The tail entries are orders, in the senses made precise in
Propositions~\ref{prop:moments-tails}, \ref{prop:matern-density} and~\ref{prop:student-t}; for
the stable row Proposition~\ref{prop:stable-family} gives the count of moments, and the order
of the density is the classical
$p(x) \sim \pi^{-1}\Gamma(\alpha+1)\sin(\pi\alpha/2)\,|x|^{-\alpha-1}$ of
\sscite{sato1999levy}, Remark~14.18, pp.~87--88, which no statement here uses. Boundedness at the origin is
Corollary~\ref{cor:origin-boundedness} for the \Matern{} and unit-step rows,
Corollary~\ref{cor:origin-smooth} for the stable row, and the explicit density of
Proposition~\ref{prop:student-t}(1) for the Student-t row.}
\label{tab:families}
\footnotesize
\setlength{\tabcolsep}{4pt}
\begin{tabular}{@{}llllll@{}}
\toprule
family & profile $k(x)$ & finite moments & tail & bounded at $0$ & lag sections \\
\midrule
Gaussian & $0$, with $a > 0$ & all & $e^{-x^2/4at^2}$ & yes & a limit only \\
\Matern{}, index $\gamma$ & $2\gamma e^{-x}$ & all & $e^{-|x|/t}$ & iff $\gamma > \tfrac12$ & iff $\gamma$ is an integer \\
Student-t, shape $a$ & a Gaussian mixture & $n < 2a$ & $|x|^{-2a-1}$ & yes & none, not a limit \\
stable, index $\alpha < 2$ & $C_\alpha^{-1}x^{-\alpha}$ & $n < \alpha$ & $|x|^{-1-\alpha}$ & yes & none, not a limit \\
unit step, cutoff $\tau$ & $\mathbf 1_{(0,\tau)}$ & all & $e^{-|x|\log|x|/\tau t}$ & no & none, not a limit \\
\bottomrule
\end{tabular}
\end{table}

% Figure: the three corner families by parameter (2026-09-18, the author's request at the
% review); the image is figures/fig-corners.png from scripts/make-fig-corners.py (numpy, scipy's
% modified Bessel function, matplotlib; closed forms except the stable densities, which are
% numerical cosine inversions checked against the Cauchy and Gaussian closed forms).
\begin{figure}[t]
\centering
\includegraphics[width=\textwidth]{fig-corners.png}
\caption{The three corner families at canonical scale $1$: the kernels (top) and the
exponents $F = -\log\hat\phi_1$ (bottom). Left, the \Matern{} family for four values of the
index $\gamma$: the Bessel $K_0$ at $\gamma = \tfrac12$, logarithmically singular at the
origin, the Laplace kernel at $\gamma = 1$, and smoother kernels beyond, approaching a
Gaussian shape; all have exponential tails, and the exponent grows like a logarithm
(Propositions~\ref{prop:matern-exponent} and~\ref{prop:matern-density}). Middle, the Student-t
family for four values of the degrees of freedom $2a$, in the scale of the proposition, in
which a larger $a$ is also a narrower kernel: the Cauchy kernel at $2a = 1$; the
tails are polynomial, of order $|x|^{-2a-1}$, and the exponent grows linearly
(Proposition~\ref{prop:student-t}). Right, the symmetric stable family for four values of the
index $\alpha$: the Cauchy kernel again at $\alpha = 1$, the Gaussian at $\alpha = 2$; the
exponent is the power $|\omega|^\alpha$ and, for $\alpha < 2$, the tails are of order
$|x|^{-1-\alpha}$ (Proposition~\ref{prop:stable-family}). The parameters vary at a fixed
canonical scale and the bandwidths are not matched, so the panels show the shapes of the
families and are not a comparison of smoothings; Figure~\ref{fig:experiment-b} compares members
at equal bandwidth.}
\label{fig:corners}
\end{figure}

\subsection{Moments and tails}\label{sec:moments-tails}

Two questions one asks of any smoothing kernel are how far it spreads and how fast it falls
off. Does it have a variance, and how many moments? Does a distant feature leak into the
measurement at an exponential or at a polynomial rate? For an admissible kernel both answers
are in the tail of the profile, and the reason is the displacement dictionary of
Remark~\ref{rem:displacement-dictionary}. The profile describes a catalogue of displacements,
$k(x)\,dx/x$ being the rate of those of size $x$; a large displacement of the kernel comes
from a large displacement in the catalogue; and the classical theory of infinitely divisible
laws makes this exact. A law has an $n$-th moment exactly when its \Levy{} measure, restricted
to the large jumps, has one (\sscite{sato1999levy}, Theorem~25.3, p.~159, in the form with a
submultiplicative weight), and when the jumps are bounded by a radius $\tau$ the tail of the
law is of the order $e^{-|x|\log|x|/\tau}$ and no lighter (\sscite{sato1999levy},
Theorem~26.1, p.~168, which states both directions). The proposition reads the two theorems
in the folded profile, adds the formula for the variance, and draws the conclusion that
matters for comparing families: a profile of bounded support gives the lightest tails the
cone allows, and they are still heavier than a Gaussian's, so the Gaussians are the only
members with Gaussian tails.

% shared with blueprint prop:moments-tails
\begin{proposition}[moments and tails]\label{prop:moments-tails}
Let $F$ be admissible with data $(a,k)$, and let $X_t$ be the displacement at canonical scale
$t$.
\begin{enumerate}
\item \emph{(Moments.)} For $n \ge 1$, $\mathbb{E}|X_t|^n < \infty$ if and only if
$\int_1^\infty x^{n-1}k(x)\,dx < \infty$. When finite, the variance is
$\mathbb{E}X_t^2 = t^2\bigl(2a + \int_0^\infty x\,k(x)\,dx\bigr)$, the two terms being the
jitter and the second moment of the catalogue.
\item \emph{(Tails.)} If $k \not\equiv 0$ and $\tau$ is the least radius beyond which $k$
vanishes, then $\mathbb{E}\exp(\alpha|X_t|\log|X_t|) < \infty$ for every $\alpha < 1/(\tau t)$
and for no $\alpha > 1/(\tau t)$: the tail is of order $e^{-|x|\log|x|/(\tau t)}$. If $k$
vanishes beyond no $\tau$, the tail is heavier than every such rate. In particular a completely monotone $k$,
being positive on all of $(0,\infty)$, puts the whole Thorin subclass on the heavy side, and no
kernel other than the Gaussians themselves has a Gaussian or lighter tail.
\end{enumerate}
\end{proposition}

The criterion of clause (1) and the two rates of clause (2) are the cited theorems read in
this paper's convention. The translation into the folded profile, the variance formula and
the last sentence of clause (2) are proved below; the variance formula uses no cited fact, the
finiteness of $\mathbb{E}X_t^2$ being a conclusion of its argument and not a hypothesis.

\begin{proof}[Proof of the parts that are this paper's own]
The \Levy{} measure of $X_t$ is the $t$-dilate of $\nu = k(x)x^{-1}dx$, folded to the two-sided
measure $\nu_2$; so
$\int_{|x|>1}|x|^n\,\nu_2(dx) = \int_{1/t}^\infty (tx)^n\,k(x)\,x^{-1}dx$, which is finite
exactly when $\int_1^\infty x^{n-1}k(x)\,dx$ is, the integrand being locally bounded away from
the origin. Applying the moment criterion with the submultiplicative weight $\max(|x|,1)^n$
gives clause (1)'s criterion: that weight agrees with $|x|^n$ on $\{|x| > 1\}$, the only region
the criterion integrates over, and has a finite integral against the law exactly when $|x|^n$
does, since $|x|^n \le \max(|x|,1)^n \le |x|^n + 1$; the weight $|x|^n$ itself is not
submultiplicative.

For the variance, what has to be evaluated is a second difference of the transform at the
origin and not a second derivative. Write $C = t^2(2a + \int_0^\infty x\,k(x)\,dx)$, finite by
hypothesis. The half-angle formula $1 - \cos u = 2\sin^2(u/2)$ gives, for every real $c$ and
every $\omega \ne 0$,
\[
  \bigl(1 - \cos c\omega\bigr)\frac{2}{\omega^2}
   = c^2\Bigl(\frac{\sin(c\omega/2)}{c\omega/2}\Bigr)^2 \le c^2 ,
\]
an identity rather than an estimate, so the same line supplies the pointwise limit $c^2$ as
$\omega \to 0$, a dominating function, and the inequality $1 - \cos u \le u^2/2$. On the
exponent side, reading it at $c = tx$ inside \ref{eq:sd-profile} and dominating by $t^2xk(x)$
gives $2\omega^{-2}F(t\omega) \to C$; hence $F(t\omega) \to 0$, and since
$ye^{-y} \le 1 - e^{-y} \le y$ for $y \ge 0$, the same limit holds for
$2\omega^{-2}(1 - e^{-F(t\omega)})$. On the measure side, $\mu_t$ being a probability measure,
$1 - \hat\mu_t(\omega) = \int(1 - \cos\omega x)\,\mu_t(dx)$; Fatou's lemma along a sequence
$\omega_m \to 0$ turns the exponent-side limit into $\int x^2\,\mu_t(dx) \le C$, so the second
moment is finite --- a conclusion of the argument, not a hypothesis of it --- and dominated
convergence with the dominating function $x^2$ then gives
$2\omega^{-2}(1 - \hat\mu_t(\omega)) \to \int x^2\,\mu_t(dx)$. The two limits are limits of the
same function, so $\mathbb{E}X_t^2 = C$. The mean vanishes because $|x| \le 1 + x^2$ makes $x$
integrable against $\mu_t$ and $\mu_t$ is reflection-symmetric.

\emph{Remark, not machine-checked.} A second route to the variance evaluates
$\mathbb{E}X_t^2 = -\partial_\omega^2 e^{-F(t\omega)}|_{\omega = 0} = t^2F''(0)$ and
differentiates \ref{eq:levy-khintchine} twice under the integral sign to obtain
$F''(0) = 2a + \int_{(0,\infty)}x^2\,\nu(dx)$. It reaches the same constant and costs a
dominating function for each of the two differentiations, where the second difference costs one
and delivers the finiteness of $\mathbb{E}X_t^2$ along the way.

For the last
sentence of (2), a completely monotone $k$ is a mixture of exponentials by Bernstein's theorem
(\S\ref{sec:thorin}): $k(x) = \int e^{-\theta x}\,V(d\theta)$ for a positive measure $V$, which is
not the zero measure because $k \not\equiv 0$. Some window $[0,\theta_0]$ then carries positive
mass, and on it $e^{-\theta x} \ge e^{-\theta_0x}$, so
$k(x) \ge e^{-\theta_0x}\,V([0,\theta_0]) > 0$ at every $x > 0$; hence
$k$ vanishes beyond no $\tau$, and the heavy branch of the tail statement applies. Finally, if
$k \equiv 0$ then $F = a\omega^2$ and the kernel is Gaussian, while if $k \not\equiv 0$ then
$F$ has a nonzero jump catalogue and the tail floor forbids a Gaussian tail: fixing $x_0 > 0$
with $k(x_0) \ne 0$, the catalogue reaches the radius $x_0$, so the floor holds at every
$\alpha' > 1/(x_0t)$, and one such $\alpha'$ suffices, since
$\alpha'|x|\log|x| \le \alpha x^2$ once $|x| \ge (2\alpha'/\alpha)^2$ --- write
$\log|x| = 2\log\sqrt{|x|} \le 2\sqrt{|x|}$ --- while the law is a probability measure, so the
divergence cannot come from the bounded part. No case distinction between a bounded and an
unbounded catalogue is needed at this step, one radius the catalogue reaches being enough.
\end{proof}

The two dials are turned by one object, the tail of $k$: polynomial moment weight by its decay,
super-exponential tail order by its support. The $\Cin$ rays and their superpositions, with
bounded support, are the best-localized members of the cone and have every moment; the price is
the echo structure of Lemma~\ref{lem:cin-delay-equation}. Completely monotone profiles give
kernels with no such echoes, computed by first-order sections when the profile is a finite
sum of exponentials with integer weights (\S\ref{sec:implementation}), at the price of tails
no lighter than exponential.

\subsection{The \texorpdfstring{\Matern{}}{Matern} family}\label{sec:matern}

The \Matern{} kernels are the functions on the line with transform
$(1 + \theta^2\omega^2)^{-\gamma}$. The name is that of the covariance functions, which are the
members with $\gamma > \tfrac12$, and this paper uses it for the whole range $\gamma > 0$. The
family has three histories. It entered statistics as the covariance functions of the
stationary random fields that solve a stochastic Helmholtz equation
(\sscite{whittle1954stationary}). It became the standard covariance family of spatial
statistics after Mat\'ern's monograph, where it is the family of correlation functions
$2(b\nu/2)^\nu K_\nu(b\nu)/\Gamma(\nu)$ in his notation (\sscite{matern1960spatial}, (2.4.7),
p.~18; the name is Stein's, who defines the class by the spectral density
$\phi(\alpha^2 + \omega^2)^{-\nu - 1/2}$ and notes that it is rational exactly at half-integer
$\nu$, \sscite{stein1999interpolation}, (14), p.~31; \sscite{guttorp2006studies} for the history; \sscite{porcu2024matern} for the present
state). And its members are the Green's functions of the operator
$(1 - \theta^2\Delta)^\gamma$, which is the link to Gaussian Markov random fields
(\sscite{lindgren2011explicit}). On the line the members of half-integer smoothness
$\nu = \gamma - \tfrac12$, which are those of integer index $\gamma$, have the spectral density
$(\lambda^2 + \omega^2)^{-\gamma}$, and its factorization into
$(\lambda + i\omega)^{-\gamma}(\lambda - i\omega)^{-\gamma}$ makes the process the output of a
linear stochastic differential equation of order $\gamma$. That is how these covariances are
computed in Gaussian process regression (\sscite{hartikainen2010kalman}, \S 4.1, (8)--(9);
\sscite{sarkka2019applied}, \S 12.3, (12.33)--(12.35)). The factorization of the member of
integer index into first-order causal and anticausal factors is therefore classical. What
Corollary~\ref{cor:matern-realization} adds is that the same factorization holds between every
two scales, a stage being a cascade of such factors with a zero at the earlier scale, and that
within the admissible class this property characterizes the family. In probability the same
functions are the densities of the symmetric variance-gamma laws, Brownian motion run for a
Gamma-distributed time, introduced as a model of asset returns (\sscite{madan1990variance};
\sscite{fischer2025variance} for the distribution theory). In image analysis they have
appeared as the Bessel potentials of the Bessel scale-space (\sscite{burgeth2005bessel}),
composed there along the order $\gamma$ at a fixed range, which is a semigroup but is not
scale invariant. What none of these sources has is the family as a scale space, composed along
the \emph{range}: with $\gamma$ fixed and the range $\theta t$ as the scale, the kernels form
an admissible family, an explicit member of the line paper
(Proposition~\ref{prop:two-members}). The stage from scale $s$ to scale $t$ has the quotient
of the two endpoint transforms as its transfer function, and it is not itself a member of the
family.

The first proposition is the computation that the rest of the subsection reads: the profile
of the family is a single exponential,
$k(x) = 2\gamma e^{-x/\theta}$, displacements of every size with an exponential cutoff at the
range, and the index $\gamma$ is their overall rate.

% shared with blueprint prop:matern-exponent
\begin{proposition}[the \Matern{} family: exponent, transforms, moments]\label{prop:matern-exponent}
For $\gamma > 0$:
\begin{enumerate}
\item \emph{(Exponent.)} $F(\omega) = \gamma\log(1 + \omega^2)$ is admissible, with Gaussian
coefficient $a = 0$ and folded profile $k(x) = 2\gamma e^{-x}$; the profile at range $\theta$ is
$k(x) = 2\gamma e^{-x/\theta}$, with exponent $\gamma\log(1 + \theta^2\omega^2)$.
\item \emph{(Transforms.)} $e^{-F(t\omega)} = (1 + t^2\omega^2)^{-\gamma}$, and the increments
have transfer functions $\bigl((1 + s^2\omega^2)/(1 + t^2\omega^2)\bigr)^{\gamma}$, rational in
$\omega$ for integer $\gamma$.
\item \emph{(Moments.)} Every moment is finite: $X_t \seq \BM_T$ with $T$ Gamma-distributed of
shape $\gamma$ and scale $2t^2$, so
$\mathbb{E}X_t^{2n} = (2n-1)!!\,(2t^2)^n\,\Gamma(\gamma+n)/\Gamma(\gamma)$, and in particular
$\mathbb{E}X_t^2 = 2\gamma t^2$.
\end{enumerate}
\end{proposition}

\begin{proof}
(1) Fix $\theta > 0$ and write $p = 1/\theta$. The integrand
$(1 - \cos\omega x)\,e^{-px}x^{-1}$ is nonnegative on $(0,\infty)$ and, since
$1 - \cos u \le u^2/2$, is bounded there by $\tfrac12\omega^2\,x\,e^{-px}$, so it is integrable
for every $\omega$; its derivative in $\omega$ is $\sin(\omega x)\,e^{-px}$, whose modulus is at
most $e^{-px}$ \emph{uniformly in} $\omega$. Differentiation under the integral sign is
therefore legitimate at every $\omega$ with one and the same dominating function, and gives
\[
  \frac{d}{d\omega}\int_0^\infty (1 - \cos\omega x)\,e^{-px}\,\frac{dx}{x}
   = \int_0^\infty \sin(\omega x)\,e^{-px}\,dx = \frac{\omega}{p^2+\omega^2},
\]
which is also the derivative of $\tfrac12\log(1 + \omega^2/p^2)$. The two functions of $\omega$
therefore differ by a constant, and both vanish at $\omega = 0$, so
\[
  \int_0^\infty (1 - \cos\omega x)\,e^{-x/\theta}\,\frac{dx}{x}
   = \tfrac12\log\bigl(1 + \theta^2\omega^2\bigr).
\]
Hence $F(\omega) = \gamma\log(1+\theta^2\omega^2)$ is of the form \ref{eq:sd-profile} with
$a = 0$ and $k(x) = 2\gamma e^{-x/\theta}$, which is nonnegative and nonincreasing on
$(0,\infty)$. Its two integrability conditions are the two dominating functions above read on
the two halves of the range: $\int_0^1 x\,k(x)\,dx \le 2\gamma$ and
$\int_1^\infty k(x)\,x^{-1}dx \le 2\gamma\int_1^\infty e^{-x/\theta}dx < \infty$. Exhibiting the
pair $(0, 2\gamma e^{-\cdot/\theta})$ is what admissibility asks for, so the argument stops
there; the range $\theta$ is carried by the computation and needs no separate dilation step.

(2) is (1) read through the similarity form of Theorem~\ref{thm:main-characterization}: the
increment transfer function is the ratio of the endpoint transforms, the denominators being
nonvanishing.

(3) Let $\Gamma_{\gamma,2t^2}$ be the Gamma law of shape $\gamma$ and scale $2t^2$, that is the
law on $(0,\infty)$ with density $u^{\gamma-1}e^{-u/2t^2}/(\Gamma(\gamma)(2t^2)^{\gamma})$. Its
Laplace transform is $\int_0^\infty e^{-\sigma u}\,\Gamma_{\gamma,2t^2}(du) =
(1 + 2t^2\sigma)^{-\gamma}$ for $\sigma \ge 0$, one Gamma integral; so the Gaussian mixture
$\int_0^\infty g_u\,\Gamma_{\gamma,2t^2}(du)$ --- the law of $\BM_T$ with $T$ so distributed ---
has cosine transform $(1 + 2t^2\cdot\omega^2/2)^{-\gamma} = (1 + t^2\omega^2)^{-\gamma}$, which
by (2) is the transform of the kernel at scale $t$. Both are symmetric probability measures, so
Proposition~\ref{prop:fourier-uniqueness} identifies them and $X_t \seq \BM_T$. Conditioning on
$T$, $\mathbb{E}X_t^{2n} = \mathbb{E}[T^n]\,\mathbb{E}[\BM_1^{2n}] = (2n-1)!!\;\mathbb{E}T^n$
with $\mathbb{E}[\BM_v^{2n}] = (2n-1)!!\,v^n$ and
$\mathbb{E}T^n = (2t^2)^n\Gamma(\gamma+n)/\Gamma(\gamma)$; the two are the same Gamma integral
read at half-integer and at integer argument. Every moment is then finite, since
$|x|^n \le 1 + x^{2n}$, and $n = 1$ with $\Gamma(\gamma+1) = \gamma\Gamma(\gamma)$ gives
$\mathbb{E}X_t^2 = 2\gamma t^2$.

\emph{Remark.} Finiteness of every moment also follows from
Proposition~\ref{prop:moments-tails}(1), the profile decaying exponentially, and the variance
from the variance formula there:
$\mathbb{E}X_t^2 = t^2\int_0^\infty x\,2\gamma e^{-x}dx = 2\gamma t^2$. That route is machine-
checked too and is independent of this one; it does not reach the general even moment.
\end{proof}

The kernel has a closed form in the modified Bessel function $K_\nu$, and it is the form under
which the family is known. The next proposition records it, with its two elementary members,
the Laplace kernel and the Bessel $K_0$, its exponential tail and its Gaussian limit. The
closed form and the tail are classical and are cited: the Fourier pair, with the
identification of the family as the symmetric variance-gamma laws, is
\sscite{fischer2025variance}, the density (1.1) and the characteristic function (2.12), with
self-decomposability in \S2.3, p.~6; the asymptotics of the modified Bessel function that give
the tail are \sscite{dlmf2026}, \S10.25(ii), (10.25.3), and
\sscite{abramowitz1964handbook}, \S9.7, 9.7.2, p.~378.

% shared with blueprint prop:matern-density
\begin{proposition}[the \Matern{} family: the kernel and its tail]\label{prop:matern-density}
For $\gamma > 0$ and $t > 0$ the kernel of the \Matern{} family at canonical scale $t$ has the
density
\[
  \phi_t(x) = \frac{1}{\sqrt\pi\,\Gamma(\gamma)\,t}
   \Bigl(\frac{|x|}{2t}\Bigr)^{\gamma - \frac12}
   K_{\gamma-\frac12}\Bigl(\frac{|x|}{t}\Bigr),
\]
for $\gamma > \tfrac12$ the \Matern{} kernel of smoothness $\gamma - \tfrac12$ and range $t$
normalised to unit mass; for every $\gamma > 0$ these are the symmetric variance-gamma laws. Here
$\gamma = 1$ gives the Laplace kernel $e^{-|x|/t}/(2t)$ at every $x \ne 0$, and
$\gamma = \tfrac12$ gives $K_0(|x|/t)/(\pi t)$ at every $x$. The origin is excluded from the first
of those two values because the display reads there as the product of $(|x|/2t)^{1/2} = 0$ with
$K_{1/2}(0)$ and does not return $(2t)^{-1}$. The difference is on a null set and leaves the
kernel, a density, unchanged.
As $|x| \to \infty$,
$\phi_t(x) \sim c_\gamma\,t^{-\gamma}|x|^{\gamma-1}e^{-|x|/t}$: exponential tails, the
heavy side of Proposition~\ref{prop:moments-tails}(2), as for every completely monotone profile. As
$\gamma \to \infty$ with $2\gamma t^2 = v$ held fixed, $\phi_t$ converges weakly to the
Gaussian of variance $v$.
\end{proposition}

The displayed function is a probability density for every
$\gamma > 0$, but a \Matern{} \emph{covariance} function is bounded at the origin, which needs
the smoothness $\gamma - \tfrac12$ to be positive, and Corollary~\ref{cor:origin-boundedness}
says the same from the side of the profile. Of the proposition's content, the two special
values, that the kernel is a density at all (Proposition~\ref{prop:kernel-regularity}) and the
Gaussian limit, from $(1 + v\omega^2/2\gamma)^{-\gamma} \to e^{-v\omega^2/2}$ and
Proposition~\ref{prop:levy-continuity}, are proved in this paper; the closed form and the tail
are the two cited facts.

The theorem says that four descriptions single out the same family. The first is the profile,
a single exponential. The second is the Thorin measure of \S\ref{sec:thorin}, a single atom:
among the mixtures of exponentials, the \Matern{} members are the pure ones. The third is the
kernel, through its transform. The fourth is the engineering one: for integer $\gamma$, every
stage of the family, from scale $s$ to scale $t$, is a rational filter with a single pole
pair, which is what lets the family be computed by recursive filters with no approximation
(\S\ref{sec:implementation}). The equivalence of the fourth with the first is the point. A
stage that is a cascade of identical first-order sections is what one would ask for on
practical grounds. Within the cone that property singles out one family, the one spatial
statistics had already chosen on grounds of its own.

% shared with blueprint thm:matern
\begin{theorem}[the \Matern{} theorem]\label{thm:matern}
For an admissible family \emph{in its canonical gauge} and $\gamma, \theta > 0$, clauses (1), (2)
and (3) are equivalent, and
when $\gamma$ is a positive integer they are equivalent to clause (4).
\begin{enumerate}
\item The profile is a single exponential, $k(x) = 2\gamma e^{-x/\theta}$, and $a = 0$.
\item The family lies in the symmetric Thorin subclass of
Proposition~\ref{prop:thorin-subclass} with
Gaussian coefficient $0$ and the one-point Thorin measure $2\gamma\,\delta_{1/\theta}$.
\item The kernels are \Matern{}: $\hat\phi_t(\omega) = (1 + \theta^2t^2\omega^2)^{-\gamma}$.
\item The stage transfer functions are $\hat\mu_{s,t}(\omega) =
\bigl((1 + \theta^2s^2\omega^2)/(1 + \theta^2t^2\omega^2)\bigr)^{\gamma}$ for all $s \le t$,
rational in $\omega$, with a single pole pair when $s < t$.
\end{enumerate}
The parameter $\theta$ is a change of gauge and may be normalised to $1$.
\end{theorem}

The equivalences hold at a given pair $(\gamma,\theta)$, and the one cited fact the proof uses
is the uniqueness of the \Levy{} pair (Proposition~\ref{prop:fourier-toolbox}(3)), which turns
an equality of exponents into an equality of profiles. Clause (4) is stated for integer
$\gamma$ only, the single-pole-pair case: at a non-integer exponent the stage transfer
function is not a ratio of polynomials. The canonical gauge is a hypothesis because clauses (3)
and (4) speak of the kernels and the stages at the parameter $t$: a \Matern{} family read in
another gauge $\chi$ has the stage transfer function with $\chi(s)$ and $\chi(t)$ in place of $s$
and $t$, and satisfies clause (1) all the same.
Proposition~\ref{prop:rational-increments} gives a class of admissible stage transfer functions;
which rational functions in general are admissible stages is open (\S\ref{sec:conclusions}), and
Remark~\ref{rem:rational-converse} shows that the first clause of that proposition has no
converse.

\begin{proof}
\emph{(1) if and only if (2).} With $k(x) = 2\gamma e^{-x/\theta}$ and $a = 0$, take
$U = 2\gamma\,\delta_{1/\theta}$. Then \ref{eq:thorin} reads
$\tfrac12\cdot 2\gamma\log(1 + \omega^2\theta^2) = \gamma\log(1+\theta^2\omega^2)$, which is the
exponent computed in Proposition~\ref{prop:matern-exponent}(1); so (1) implies (2) with that
one-point measure, the factor $2$ of the folded profile cancelling the $\tfrac12$ of
\ref{eq:thorin}. Conversely, a one-point $U = 2\gamma\,\delta_{1/\theta}$ makes the right-hand
side of \ref{eq:thorin} that same exponent, so $F(\omega) = \gamma\log(1+\theta^2\omega^2)$, and
the next paragraph's uniqueness argument returns the single exponential.

\emph{(1) implies (3)} is Proposition~\ref{prop:matern-exponent}(1)--(2) after the gauge change
$t \mapsto \theta t$, which Proposition~\ref{prop:matern-exponent}(1) supplies directly at range
$\theta$.

\emph{(3) implies (1).} Reading the transform at the canonical scale $t = 1$ and taking
logarithms turns clause (3) into $F(\omega) = \gamma\log(1+\theta^2\omega^2)$; the other scales
say nothing further. The computation in Proposition~\ref{prop:matern-exponent}(1) shows that
this exponent has the pair $(0, 2\gamma e^{-x/\theta}x^{-1}dx)$, and the pair of an exponent is
unique (Proposition~\ref{prop:fourier-toolbox}(3)), so $a = 0$ and the \Levy{} measure is that
one. Uniqueness of the pair is an identity of \emph{measures}, so it gives
$k(x) = 2\gamma e^{-x/\theta}$ for almost every $x > 0$; the pointwise identity on $(0,\infty)$
follows because $k$ is nonincreasing there while $2\gamma e^{-\cdot/\theta}$ is continuous.
Indeed the set on which the two agree is co-null, hence meets every subinterval of $(0,\infty)$;
approaching a given $x > 0$ from the left through that set gives
$k(x) \le 2\gamma e^{-y/\theta}$ for points $y < x$ arbitrarily close to $x$, whence
$k(x) \le 2\gamma e^{-x/\theta}$, and approaching from the right gives the reverse inequality.
(Monotonicity of the comparison function is not used; continuity of one side and monotonicity
of the other suffice.)

\emph{(3) if and only if (4)}, for integer $\gamma$, is Proposition~\ref{prop:matern-exponent}(2)
read in both directions. Forwards, the cascade law makes the increment transfer function the
ratio of the two endpoint transforms, the denominators being nonvanishing; at the collapsed pair
$s = t = 0$ the same law reads $\hat\mu_{0,1} = \hat\mu_{0,1}\hat\mu_{0,0}$ with
$\hat\mu_{0,1} > 0$, so $\hat\mu_{0,0} = 1$, which is clause (4) there. Backwards, clause (4) at
$s = 0$ is the endpoint transform of clause (3); no limit in $s$ is needed.
\end{proof}

The engineering reading of the theorem is the next corollary. One of its terms belongs to
\S\ref{sec:implementation}: a first-order forward--backward section is a first-order recursive
filter run forward along the signal, together with its mirror image run backward, and a cascade
of $\gamma$ identical such pairs has the transfer function
$\bigl((1+p^2\omega^2)/(1+q^2\omega^2)\bigr)^\gamma$, with $q$ the range of the pole and $p$ that
of the zero. The realization half of the corollary is proved in that section, where the sections
are constructed. What the corollary does \emph{not} prescribe is the pair: the pole and the zero
are not tied in advance to the two scales of the stage, and the third clause asks for one kernel
from the origin and nothing else. That the zero then vanishes is the content of the argument.

% shared with blueprint cor:matern-realization
\begin{corollary}[the \Matern{} families of integer index are the ones realized by identical first-order sections]
\label{cor:matern-realization}
Let $(\Phis{s}{t})$ be an admissible family in its canonical gauge and let $\gamma$ be a positive
integer. Call a symmetric kernel \emph{a cascade of $\gamma$ identical forward--backward pairs of
first-order sections} when its transform is
$\bigl((1+p^2\omega^2)/(1+q^2\omega^2)\bigr)^{\gamma}$ for some $p, q \ge 0$, the pair being the
first-order forward--backward section pair of Proposition~\ref{prop:matern-cascade} with pole
range $q$ and zero range $p$. Then the following are equivalent.
\begin{enumerate}
\item The family is \Matern{} of index $\gamma$: its exponent is
$F(\omega) = \gamma\log(1+\theta^2\omega^2)$ for some $\theta > 0$.
\item Every stage is such a cascade: for all $0 \le s < t$ the increment $\mu_{s,t}$ is a cascade
of $\gamma$ identical pairs, with $p$ and $q$ depending on the pair of scales.
\item One kernel from the origin is such a cascade: for some $t > 0$ the kernel $\phi_t$ is.
\end{enumerate}
In (1) the pair of the stage from $s$ to $t$ is $(p,q) = (\theta s, \theta t)$; in (3) the zero
necessarily vanishes, $p = 0$ and $q = \theta t$. Such a family has every moment finite and
exponential tails.
\end{corollary}

The word \emph{identical} is what carries the corollary, and it cannot be dropped. An admissible
family whose kernel from the origin is a cascade of first-order forward--backward pairs of
different ranges need not be \Matern{}: the second example of
Remark~\ref{rem:rational-converse}, of exponent $2\log(1+\omega^2) - \log(1+\omega^2/4)$, has the
kernel transform $\bigl[(1+t^2\omega^2/4)/(1+t^2\omega^2)\bigr]\cdot(1+t^2\omega^2)^{-1}$ from the
origin at scale $t$, a cascade of two pairs, one carrying a zero of range $t/2$ and one a pure
pole of range $t$. The same happens with three pairs at $k(x) = 6e^{-x} - 2e^{-2x}$, which is
positive ($6e^{-x} \ge 2e^{-x} > 2e^{-2x}$) and nonincreasing
($k'(x) = -2e^{-x}(3-2e^{-x}) < 0$) and has
$F(\omega) = 3\log(1+\omega^2) - \log(1+\omega^2/4)$ by
Proposition~\ref{prop:matern-exponent}(1) at the two rates.

\begin{proof}
\emph{(1) implies (2).} A \Matern{} family of positive integer index $\gamma$ and range $\theta$
has, by Proposition~\ref{prop:matern-cascade} read in the gauge $t \mapsto \theta t$, the stage
transfer functions $\bigl((1+\theta^2s^2\omega^2)/(1+\theta^2t^2\omega^2)\bigr)^\gamma$ at every
pair of scales, realized by $\gamma$ identical forward--backward pairs; so every stage is a
cascade of the kind named, with $(p,q) = (\theta s,\theta t)$ and $p < q$.

\emph{(2) implies (3)} is the stage from $s = 0$, which in the canonical gauge is the kernel
$\phi_t = \mu_{0,t}$ from the origin, at any $t > 0$.

\emph{(3) implies (1).} Suppose
$\hat\phi_t(\omega) = \bigl((1+p^2\omega^2)/(1+q^2\omega^2)\bigr)^\gamma$ for some $t > 0$ and
some $p,q \ge 0$. First $q > p$: the transform of a probability measure is bounded by $1$ in
modulus, so the ratio $(1+p^2\omega^2)/(1+q^2\omega^2)$ is bounded on the line, which forces
$q \ge p$, and $q = p$ would make $\hat\phi_t \equiv 1$, hence $\phi_t = \delta_0$ and
$\Phis{0}{t} = \mathrm{Id}$, against (ND). Next $p = 0$: by
Proposition~\ref{prop:kernel-regularity} the kernel from the origin is absolutely continuous, so
its transform tends to $0$ at infinity --- the Fourier transform of an integrable function does
--- while
\[
  \Bigl(\frac{1+p^2\omega^2}{1+q^2\omega^2}\Bigr)^{\gamma}
   \ \longrightarrow\ \Bigl(\frac pq\Bigr)^{2\gamma} \qquad (|\omega| \to \infty),
\]
a positive number unless $p = 0$. So $p = 0$ and $q > 0$, and the canonical gauge reads
$e^{-F(t\omega)} = (1+q^2\omega^2)^{-\gamma}$, that is
$F(t\omega) = \gamma\log(1+q^2\omega^2)$; substituting $\omega \mapsto \omega/t$ gives clause (1)
with $\theta = q/t$.

\emph{The index is not free.} If the kernels from the origin at two scales are cascades of
$\gamma$ and $\gamma'$ identical pairs, the paragraph just proved gives
$\gamma\log(1+\theta^2\omega^2) = \gamma'\log(1+\theta'^2\omega^2)$; dividing by $\log|\omega|$
and letting $|\omega| \to \infty$ gives $\gamma = \gamma'$. So the number of pairs is the index,
at every scale.

\emph{The consequences.} Every moment is finite by
Proposition~\ref{prop:matern-exponent}(3), and the tail is exponential by
Proposition~\ref{prop:matern-density}.
\end{proof}

% blueprint rem:matern-literature, copied; moved here from the end of the section and given its
% citekeys paper-side (2026-09-18, the author's review: each corner placed in its literature).
\begin{remark}[the literature of the \Matern{} corner]\label{rem:matern-literature}
The kernels of Theorem~\ref{thm:matern} are the \Matern{} covariance functions of spatial
statistics, defined there by the spectral density $(1+\theta^2\omega^2)^{-\gamma}$ or as the
Green's functions of a fractional Helmholtz operator (\sscite{whittle1954stationary};
\sscite{lindgren2011explicit}; \sscite{porcu2024matern}). In image analysis the same kernels
appear as Bessel potentials, and the Bessel scale-space built from them
(\sscite{burgeth2005bessel}) composes the family along the \emph{order} $\gamma$ at fixed
range --- an additive semigroup in $\gamma$ that is explicitly not scale invariant --- with no
composition law stated for the range family at fixed order. In natural-image statistics they
appear as Bessel K forms, a statistical model of filter responses rather than a kernel family
(see \sscite{pedersen2005alpha}). In none of these is the family derived from axioms on
measurement, and the hemigroup is what makes the derivation possible: composed along the range
the family is a two-parameter cascade whose increments are not stationary in the scale
parameter, so no semigroup axiomatics can reach it, which is why the literature composed along
the order instead. The probabilistic identity, the symmetric variance-gamma laws
(\sscite{madan1990variance}), is what the classification actually sees.
\end{remark}

\subsection{The symmetric Thorin subclass}\label{sec:thorin}

A function on $(0,\infty)$ is \emph{completely monotone} when it is smooth and its derivatives
alternate in sign, $(-1)^n k^{(n)} \ge 0$ for every $n$. By Bernstein's theorem, cited below,
these are the mixtures of exponentials, $k(x) = \int e^{-\theta x}\,U(d\theta)$ with a positive
measure $U$. Each exponential is a
\Matern{} profile, so these are the superpositions of \Matern{} members over their ranges, and
the mixing measure $U$ is called the \emph{Thorin measure}. The class is classical. On the
half-line it is Thorin's class of generalized gamma convolutions, introduced to prove the
infinite divisibility of the Pareto and the lognormal laws (\sscite{thorin1977pareto};
\sscite{thorin1977lognormal}) and developed into a theory by Bondesson
(\sscite{bondesson1992generalized}); the class here is its symmetric extension to the line.

Three things make it the organising object of this section. All three named families lie in
it, the \Matern{} family as its atoms and the Student-t family under the hypothesis of
Remark~\ref{rem:student-ggc}. It is the part of the cone on which the bridge of
\S\ref{sec:bridge} is a bijection: a causal family with a completely monotone delay profile and
its image determine each other, with an explicit map between their Thorin measures. And it
contains the families realized by first-order sections, as the members whose Thorin measure is a finite
sum of atoms of even integer mass (Proposition~\ref{prop:thorin-machine}). The
proposition gives the representation of the exponent that goes with the mixture,
\ref{eq:thorin}, in which every atom of $U$ contributes one \Matern{} exponent; clause (3) is
the bijection, and clause (4) shows that the subclass is strictly smaller than the image of the
bridge, so that the chain from the Thorin subclass through the subordinated slice to the cone
is strict at both steps. With the \Matern{} members at its start, the chain and the members
that make each step strict are
\[
\begin{array}{ccccccc}
\text{\Matern{} members} & \subsetneq & \text{Thorin subclass} & \subsetneq & \mathcal{S}
  & \subsetneq & \text{admissible cone} \\[2pt]
 & & \uparrow & & \uparrow & & \uparrow \\[2pt]
 & & \text{stable} & & k(x) = 2\,\erfc\bigl(x/\sqrt{2\tau}\bigr) & & \Cin\ \text{rays}
\end{array}
\]
where each member named in the lower row lies in the class above it and not in the class to
the left of that one. The Student-t members sit beside the stable ones under the hypothesis of
Remark~\ref{rem:student-ggc}, and unconditionally nowhere in this chain. The first step is strict
because the \Matern{} members are the atoms of
the subclass, the second by clause (4), and the third by
Proposition~\ref{prop:bridge-strictness}(3). Figure~\ref{fig:filtration} divides the last step
further, by dimension. The one cited fact is Bernstein's theorem, that a completely monotone
function is the Laplace transform of a positive measure (\sscite{schilling2012bernstein},
Theorem~1.4, p.~3; \sscite{feller2009introduction}, Vol.~2, \S XIII.4, Theorems~1 and~1a,
pp.~439--440).

% shared with blueprint prop:thorin-subclass
\begin{proposition}[the symmetric Thorin subclass]\label{prop:thorin-subclass}
For an admissible $F$ with data $(a,k)$ the following are equivalent.
\begin{enumerate}
\item $k$ is completely monotone.
\item $F$ has a symmetric Thorin representation
\begin{equation}\label{eq:thorin}
  F(\omega) = a\,\omega^2
   + \tfrac12\int_{(0,\infty)} \log\Bigl(1 + \frac{\omega^2}{\theta^2}\Bigr)\,U(d\theta),
  \qquad U \ge 0, \quad
  \int_{(0,\infty)}\log\Bigl(1 + \frac{1}{\theta^2}\Bigr)\,U(d\theta) < \infty .
\end{equation}
\end{enumerate}
A third member of the equivalence --- that the kernels are symmetric extended generalized gamma
convolutions in Thorin's sense --- is stated apart, as Lemma~\ref{lem:thorin-ggc}, because it
rests on the reading of Thorin's definition rather than on Bernstein's theorem.

Moreover:
\begin{enumerate}
\item[(3)] \emph{(The bridge on Thorin subclasses.)} If $F_{\mathrm I}$ is causally admissible
with completely monotone profile $k_{\mathrm I} = \int e^{-\theta u}\,U_{\mathrm I}(d\theta)$,
then $F_{\mathrm I}(\cdot^2/2)$ is a symmetric Thorin member with $a = \tfrac12 b_0$ and $U$
equal to \emph{twice} the image of $U_{\mathrm I}$ under $\theta \mapsto \sqrt{2\theta}$; the
map $U_{\mathrm I} \mapsto U$ is linear and injective, and carries the causal Thorin measures
onto the Thorin measures of the symmetric Thorin members, so the subordination bridge restricts
to a linear bijection of the causal Thorin subclass onto the symmetric one.
\item[(4)] \emph{(Strictness the other way.)} The subordinated slice $\mathcal{S}$ of
Proposition~\ref{prop:bridge-strictness} is strictly larger than the symmetric Thorin subclass:
the image of the causal ray with profile $k_{\mathrm I} = \mathbf 1_{(0,\tau)}$ has the explicit
profile $k(x) = 2\,\erfc\bigl(x/\sqrt{2\tau}\bigr)$, which is not completely monotone. So the
symmetric Thorin subclass is strictly inside $\mathcal{S}$, which is strictly inside the
admissible cone.
\end{enumerate}
\end{proposition}

Clause (3) is classical. Both directions, with the correspondence of the measures, are
Bondesson's Theorem~7.3.1, p.~115: if $\varphi$ is the moment generating function of a
generalized gamma convolution then $\psi(s) = \varphi(s^2/2)$ is that of a symmetric extended
generalized gamma convolution, and conversely; and the proof on pp.~115--116 computes the
spectral measure of $\psi$ as the image of that of $\varphi$ under $t \mapsto \pm\sqrt{2t}$, the
left extremity of $\varphi$ becoming half the coefficient of $s^2$ in $\psi$. Read at imaginary
argument --- which is Lemma~\ref{lem:thorin-ggc} --- that is the clause. What belongs to this
paper is the folded normalisation, the linearity and injectivity of the measure map, and the
verification that the constructed causal datum is causally admissible.

What is taken from Bernstein's theorem is its existence half. The uniqueness of the Thorin
measure is not taken from it: it is Proposition~\ref{prop:laplace-uniqueness-locally-finite}.
The computation that turns the mixture into \ref{eq:thorin}, the correspondence between the
two integrability conditions, and clauses (3) and (4), which are computations on the bridge
profile \ref{eq:bridge-profile}, are proved below. An atom of $U$ at $\theta = 0$ is excluded
because it would make $k$ bounded below and $\int_1^\infty k(x)x^{-1}dx$ infinite.

\begin{proof}[Proof of the parts that are this paper's own]
\emph{(1) implies (2).} By Bernstein's theorem, $k(x) = \int_{[0,\infty)} e^{-\theta x}\,U(d\theta)$
for a positive measure $U$ on $[0,\infty)$; an atom at $\theta = 0$ would make $k$ bounded below
by a positive constant, so that $\int_1^\infty k(x)x^{-1}dx = \infty$, against admissibility. So
$U$ lives on $(0,\infty)$, and it is $\sigma$-finite, its transform being finite at $x = 1$.
Tonelli with the computation of Proposition~\ref{prop:matern-exponent}(1) gives
\[
  \int_0^\infty (1-\cos\omega x)\,k(x)\,\frac{dx}{x}
  = \int_{(0,\infty)}\Bigl(\int_0^\infty (1-\cos\omega x)\,e^{-\theta x}\,\frac{dx}{x}\Bigr) U(d\theta)
  = \tfrac12\int_{(0,\infty)}\log\Bigl(1 + \frac{\omega^2}{\theta^2}\Bigr)U(d\theta),
\]
which is \ref{eq:thorin}. The integrability condition of \ref{eq:thorin} is not a further
hypothesis but a consequence of that identity: read at $\omega = 1$ the identity says that the
integral is $2(F(1) - a)$, and $F(1)$ is finite for every admissible $F$ by
Lemma~\ref{lem:quadratic-growth}. Three regimes read off from it what the classical form of the
condition states separately: $\log(1+\theta^{-2}) \ge 2\log(1/\theta)$ near $0$,
$\ge \tfrac12\theta^{-2}$ at infinity, and $\ge \log 2$ on $(0,1]$, so that
$\int_{(0,1]}\log(1/\theta)\,U(d\theta)$ and $\int_{(1,\infty)}\theta^{-2}\,U(d\theta)$ are finite
and $U$ is finite on $(0,1]$, hence on every compact subset of $(0,\infty)$.

\emph{(2) implies (1).} Given \ref{eq:thorin} with $U$ carried by $(0,\infty)$, set
$k'(x) := \int_{(0,\infty)} e^{-\theta x}\,U(d\theta)$, which is finite for every $x > 0$: on
$(0,1]$ the integrand is at most $1$ and $U$ is finite there, and on $(1,\infty)$ one has
$e^{-\theta x} \le C_x\theta^{-2}$ against the second integral, both facts being the three regimes
of the previous paragraph. It is the finiteness of $U$ on $(0,1]$ that is doing the work here, and
it is why the displayed condition is the one printed: the two integrals alone leave $U$ free to
put infinite mass just below $\theta = 1$. The displayed Tonelli, read backwards, says that $k$ and $k'$ have the same
compensated integral $\int_0^\infty(1-\cos\omega x)\,\cdot\,dx/x$ at every frequency. That $k'$
is the profile of a symmetric \Levy{} pair at all --- that is, that
$\int_0^\infty (1\wedge x^2)\,k'(x)\,dx/x < \infty$ --- follows by averaging that identity over
$\omega$ against a standard Gaussian law: Tonelli replaces the weight $1 - \cos\omega x$ by
$1 - e^{-x^2/2}$, which lies between $\tfrac13(1 \wedge x^2)$ and $1 \wedge x^2$, so the \Levy{}
condition transfers from $k$ to $k'$. Uniqueness of the pair
(Proposition~\ref{prop:fourier-toolbox}(3)) then identifies the two \Levy{} measures, hence $k$
and $k'$ almost everywhere on $(0,\infty)$, and hence everywhere there, $k$ being nonincreasing
and $k'$ continuous. So $k$ is a Laplace transform, and Bernstein's theorem read the other way
returns complete monotonicity.

\emph{Uniqueness of $U$.} If $U'$ on $(0,\infty)$ also satisfies \ref{eq:thorin}, the previous
paragraph applied to $U'$ gives $k(x) = \int e^{-\theta x}\,U'(d\theta)$ as well, so $U$ and $U'$
have the same Laplace transform on $(0,\infty)$ and it is finite there; by
Proposition~\ref{prop:laplace-uniqueness-locally-finite}, $U' = U$. This does not spend the
uniqueness half of Bernstein's theorem: the Laplace uniqueness the argument needs is proved in
the line paper.

\emph{(3).} With $k_{\mathrm I}(u) = \int e^{-\theta u}\,U_{\mathrm I}(d\theta)$,
\ref{eq:bridge-profile} and Tonelli give
\[
  k(x) = \int U_{\mathrm I}(d\theta)\;2x\int_0^\infty g_u(x)\,e^{-\theta u}\,\frac{du}{u}
       = 2\int e^{-x\sqrt{2\theta}}\,U_{\mathrm I}(d\theta),
\]
by the first-passage identity used in Proposition~\ref{prop:bridge-families}(2); substituting
$\theta' = \sqrt{2\theta}$ this is $\int e^{-\theta'x}\,U(d\theta')$ with $U$ twice the image
of $U_{\mathrm I}$, so $k$ is completely monotone, and the Gaussian coefficient is
$\tfrac12 b_0$ by Lemma~\ref{lem:bridge-exponents}. That $U_{\mathrm I}$ is $\sigma$-finite,
which the Tonelli needs, is not a further hypothesis: $k_{\mathrm I}(1)$ is a real number, so
the transform of $U_{\mathrm I}$ is finite at $u = 1$ against a weight vanishing nowhere.

The identification of $k$ deserves a sentence, because the exponent determines it only up to a
null set. What Proposition~\ref{prop:fourier-toolbox}(3) returns, applied to the subordinated
exponent and to any admissible $F$ with that exponent, is equality of the two \Levy{} measures,
hence of the two profiles almost everywhere on $(0,\infty)$; the conclusions of this clause ---
complete monotonicity, which is a statement about every derivative at every point --- are
pointwise. The passage is the one Theorem~\ref{thm:matern} makes: the two profiles are
nonincreasing and the Laplace transform $\int e^{-\theta'x}U(d\theta')$ is continuous on
$(0,\infty)$ by dominated convergence, and a nonincreasing function almost everywhere equal to
a continuous one equals it everywhere, the good set meeting every subinterval on each side of a
given point.

The factor $2$ is the folding convention and is checked on the \Matern{} member: the causal
Gamma data are $U_{\mathrm I} = \gamma\delta_1$, whence $U = 2\gamma\delta_{\sqrt2}$ and
\ref{eq:thorin} returns
$\tfrac12\cdot 2\gamma\log(1+\omega^2/2) = \gamma\log(1+\omega^2/2)$, which is
Proposition~\ref{prop:bridge-families}(2); the image without the factor would return half of it.

\emph{The bridge is onto the symmetric Thorin subclass.} Let $F$ be admissible with data $(a,k)$
and Thorin measure $U$ on $(0,\infty)$ as in \ref{eq:thorin}. Put $b_0 := 2a$, let
$U_{\mathrm I}$ be half the image of $U$ under $\theta' \mapsto \theta'^2/2$, and let
$k_{\mathrm I}(u) := \int e^{-\theta u}\,U_{\mathrm I}(d\theta)$ for $u > 0$. The transform is
finite: the exponents satisfy
$u\theta'^2/2 - (u\theta'/2 - u/2) = \tfrac u2\bigl((\theta' - \tfrac12)^2 + \tfrac34\bigr) \ge 0$,
so $k_{\mathrm I}(u) \le \tfrac12 e^{u/2}\int e^{-u\theta'/2}\,U(d\theta')$, which the three
regimes of (1) implies (2) make finite. Hence $k_{\mathrm I}$ is nonnegative and nonincreasing on
$(0,\infty)$, and $U_{\mathrm I}$ is $\sigma$-finite. The exponent of $(b_0, k_{\mathrm I})$ is
computed rather than identified. For $u > 0$ and $\sigma \ge 0$ one has
$(1 - e^{-\sigma u})/u = \int_0^\sigma e^{-su}\,ds$, and integrating out $u$ first,
\[
  \int_0^\infty (1 - e^{-\sigma u})\,k_{\mathrm I}(u)\,\frac{du}{u}
  = \int U_{\mathrm I}(d\theta)\int_0^\sigma \frac{ds}{s + \theta}
  = \int \log\Bigl(1 + \frac{\sigma}{\theta}\Bigr)\,U_{\mathrm I}(d\theta)
\]
by Tonelli, twice. At $\sigma = \omega^2/2$ the substitution $\theta = \theta'^2/2$ turns
$\log(1 + \sigma/\theta)$ into $\log(1 + \omega^2/\theta'^2)$ and $U_{\mathrm I}$ into
$\tfrac12 U$, and $b_0\sigma = a\omega^2$, so $F_{\mathrm I}(\omega^2/2)$ is the right-hand side
of \ref{eq:thorin}, which is $F(\omega)$. At $\sigma = 1$ the display gives admissibility. The two
integrability conditions of the causal cone together say
$\int_0^\infty \min(1,u)\,k_{\mathrm I}(u)\,du/u < \infty$, and $\min(1,u) \le e\,(1 - e^{-u})$
for $u \ge 0$: for $u \le 1$ because $e\,(1 - e^{-u}) \ge e^u(1 - e^{-u}) = e^u - 1 \ge u$, and
for $u \ge 1$ because $e\,(1 - e^{-u}) \ge e - 1 \ge 1$. So the conditions are bounded by
$e\int \log(1 + 1/\theta)\,U_{\mathrm I}(d\theta)
= \tfrac e2\int \log(1 + 2/\theta'^2)\,U(d\theta')$, and $\log(1 + 2x) \le 2\log(1 + x)$, which
is $1 + 2x \le (1 + x)^2$, bounds that by $e$ times the Thorin integral of \ref{eq:thorin} at
$\omega = 1$, finite because $F(1)$ is. Finally $\sqrt{2\cdot\theta'^2/2} = \theta'$ on
$(0,\infty)$, so $U$ is twice the image of $U_{\mathrm I}$ under $\theta \mapsto \sqrt{2\theta}$.
So every symmetric Thorin member is the bridge image of a causal Thorin member, with Gaussian
coefficient and Thorin measures related as in (3); no uniqueness theorem and no part of
Bernstein's theorem is used.

\emph{The measure map is a linear bijection.} It is linear because taking an image measure and
multiplying by $2$ both are. It is injective: if twice the images of $U_{\mathrm I}$ and
$U_{\mathrm I}'$ under $\theta \mapsto \sqrt{2\theta}$ agree, so do the images, and the image of
either under $\theta' \mapsto \theta'^2/2$, which inverts $\theta \mapsto \sqrt{2\theta}$ on
$(0,\infty)$, returns the measure itself. Surjectivity onto the symmetric Thorin measures is the
previous paragraph.

\emph{A second route to the correspondence.} The two conditions also correspond directly, without
passing through an exponent. Write
$C_{\mathrm I} := \int\log(1 + \theta^{-1})\,U_{\mathrm I}(d\theta)$ for the causal condition and
$C := \int\log(1 + \theta'^{-2})\,U(d\theta')$ for \ref{eq:thorin}'s. Since $U$ is twice the image
of $U_{\mathrm I}$ under $\theta \mapsto \sqrt{2\theta}$, the substitution $\theta'^2 = 2\theta$
gives $C = 2\int\log(1 + v)\,U_{\mathrm I}(d\theta)$ with $v := 1/2\theta$, while
$C_{\mathrm I} = \int\log(1 + 2v)\,U_{\mathrm I}(d\theta)$; and
$\log(1+v) \le \log(1+2v) \le 2\log(1+v)$, the second inequality being $1 + 2v \le (1+v)^2$. So
$C_{\mathrm I} \le C \le 2C_{\mathrm I}$, and each condition holds exactly when the other does.
The route above does not compare the conditions at all: it reads both off the finiteness of one
exponent value.

\emph{(4).} The causal ray has $k_{\mathrm I} = \mathbf 1_{(0,\tau)}$, and
\ref{eq:bridge-profile} gives $k(x) = 2x\int_0^\tau g_u(x)\,du/u$. Substituting $v = x^2/2u$,
so that $du/u = -dv/v$ and $g_u(x) = \sqrt{v/\pi}\,x^{-1}e^{-v}$,
\[
  \int_0^\tau g_u(x)\,\frac{du}{u}
  = \frac{1}{\sqrt\pi\,x}\int_{x^2/2\tau}^\infty v^{-1/2}e^{-v}\,dv
  = \frac{1}{x}\,\erfc\Bigl(\frac{x}{\sqrt{2\tau}}\Bigr),
\]
so $k(x) = 2\,\erfc(x/\sqrt{2\tau})$, identified pointwise by the passage of clause (3) ---
$\erfc$ is continuous, being an integral over a moving half-line of an integrable function.

It is not completely monotone, and one bound settles it. Suppose it were; then by Bernstein's
theorem it is $\int e^{-\theta x}\,V(d\theta)$ for a positive $V$ on $[0,\infty)$, and $V \ne 0$
because $\erfc > 0$. Some window carries mass: $V(-\infty, n] > 0$ for an integer $n$, and that
mass is finite because the transform is. On that window $e^{-\theta x} \ge e^{-nx}$, so
\[
  k(x) \;\ge\; c\,e^{-nx}, \qquad c := V(-\infty,n] > 0,
\]
for every $x > 0$. Against that, for $y \ge 1$ the comparison $e^{-v^2} \le e^{-yv}$ on
$(y,\infty)$ gives $\int_y^\infty e^{-v^2}dv \le e^{-y^2}/y \le e^{-y^2}$, so
$\erfc y \le 2e^{-y^2}$ and $k(x) \le 4e^{-x^2/2\tau}$ once $x \ge \sqrt{2\tau}$. Take
$x \ge \max\bigl(\sqrt{2\tau},\,2\tau(n+1),\,4/c + 1\bigr)$. The middle bound gives
$x^2/2\tau \ge (n+1)x$ and the last gives $4e^{-x} \le 4/x < c$, whence
\[
  k(x) \;\le\; 4e^{-x^2/2\tau} \;\le\; 4e^{-x}\,e^{-nx} \;<\; c\,e^{-nx} \;\le\; k(x),
\]
which is absurd. Its image is admissible by Lemma~\ref{lem:bridge-exponents}, so it lies in
$\mathcal{S}$ and not in the Thorin subclass; the other strict inclusion is
Proposition~\ref{prop:bridge-strictness}(4).

\emph{Remark, not machine-checked.} The direct route reads the same conclusion off the
asymptotic $\erfc(x/\sqrt{2\tau}) \sim \sqrt{2\tau/\pi}\,x^{-1}e^{-x^2/2\tau}$, which says more
than the contradiction needs: an upper bound of Gaussian type at one sufficiently large $x$ is
enough, and the bound above costs a monotone comparison where the asymptotic costs Watson's
lemma.
\end{proof}

Where the named families sit in the chain of (4) is said by their own statements. The
\Matern{} family is the one-atom case of (2), so it lies in the symmetric Thorin subclass. The
Thorin density of the stable family is Proposition~\ref{prop:stable-family}. The membership of
the Student-t family is Proposition~\ref{prop:student-t}(3), under the hypothesis stated
there. The $\Cin$ rays lie outside $\mathcal{S}$ by Proposition~\ref{prop:bridge-strictness}(4).

The bijection in clause (3) is why the causal corners of \sscite{fagerstrom2026hemigroup} land
on named spatial families with the same organisation. The causal Gamma family is the single
atom of the causal Thorin subclass, and the \Matern{} family is the single atom of the
symmetric one. The stable families correspond under $\theta \mapsto \sqrt{2\theta}$. The causal
Bessel family maps to a Student-t family that is Thorin too, under the hypothesis of
Remark~\ref{rem:student-ggc}. Clause (4) adds
that the bridge is a bijection of the \emph{Thorin} subclasses and only an injection of the
whole cones.

It remains to connect \ref{eq:thorin} with Thorin's own definition, so that the name is
earned. Bondesson defines the extended generalized gamma convolutions by a two-sided
exponent: a moment generating function with a measure $V$ on the whole line and a compensating
term for each $\theta$ (\sscite{bondesson1992generalized}, Chapter~7, the definition on p.~105
with (7.1.1) and the integrability condition (7.1.2); the generalized gamma convolutions
themselves are Theorem~3.1.1, \S3.1, p.~30). For a symmetric law the contributions of $\theta$
and $-\theta$ pair up, the compensating terms cancel, and what is left at imaginary argument
is the right-hand side of \ref{eq:thorin} with $U$ the fold of $V$. The first lemma is that
computation, and it is elementary.

% shared with blueprint lem:thorin-two-sided
\begin{lemma}[Thorin's two-sided exponent at $\sigma = i\omega$]\label{lem:thorin-two-sided}
Let $V \ge 0$ be a measure on $\RR$ with $V\{0\} = 0$, invariant under $\theta \mapsto -\theta$,
and let $U := 2\,V|_{(0,\infty)}$ be its fold onto the positive half-line. For $\theta \ne 0$ and
$\omega \in \RR$,
\[
  \frac{\theta}{\theta - i\omega}\cdot\frac{-\theta}{-\theta - i\omega}
   = \Bigl(1 + \frac{\omega^2}{\theta^2}\Bigr)^{-1},
  \qquad
  -\log\Bigl|\frac{\theta}{\theta - i\omega}\Bigr|
   = \frac12\log\Bigl(1 + \frac{\omega^2}{\theta^2}\Bigr),
  \qquad
  \frac{i\omega\,\theta}{1+\theta^2} + \frac{i\omega\,(-\theta)}{1+(-\theta)^2} = 0 ,
\]
and consequently
\[
  -\int_{\RR}\log\Bigl|\frac{\theta}{\theta - i\omega}\Bigr|\,V(d\theta)
   = \frac12\int_{(0,\infty)}\log\Bigl(1 + \frac{\omega^2}{\theta^2}\Bigr)\,U(d\theta) .
\]
So Thorin's two-sided exponent, read at $\sigma = i\omega$ with $b = 0$ and $V$ symmetric, is the
right-hand side of \ref{eq:thorin} at $(a, U)$ with $a = \tfrac12c$: the first identity pairs
$\pm\theta$ without choosing a branch of the logarithm, the third cancels the compensating terms,
and the second says that the real part alone already supplies the folded integrand.
\end{lemma}

\begin{proof}
The three identities are pointwise computations at a fixed $\theta \ne 0$ and $\omega \in \RR$.
For the first, the two denominators multiply to $(\theta - i\omega)(-\theta - i\omega)
= -(\theta - i\omega)(\theta + i\omega) = -(\theta^2 + \omega^2)$, and the two numerators to
$-\theta^2$, so the product is $\theta^2/(\theta^2 + \omega^2) = (1 + \omega^2/\theta^2)^{-1}$,
a positive real number: no branch of a logarithm is chosen because none is taken. For the
second, $|\theta - i\omega| = \sqrt{\theta^2 + \omega^2}$, so
$-\log|\theta/(\theta - i\omega)| = \tfrac12\log\bigl((\theta^2+\omega^2)/\theta^2\bigr)
= \tfrac12\log(1 + \omega^2/\theta^2)$. The third is $\theta + (-\theta) = 0$ over the common
denominator $1 + \theta^2$, which is even in $\theta$.

For the fold, the integrand of the left-hand side is, by the second identity,
$\tfrac12\log(1 + \omega^2/\theta^2)$, an even function of $\theta$ that is finite off the
origin; $V$ gives the origin no mass and is invariant under $\theta \mapsto -\theta$, so the
integral over $\RR$ is twice the integral over $(0,\infty)$, and twice $V|_{(0,\infty)}$ is $U$.
That is the displayed identity. Read at Thorin's exponent with $b = 0$, the compensating terms
cancel in the pairs $\pm\theta$ by the third identity, the Gaussian term $\tfrac12 c\,\sigma^2$
at $\sigma = i\omega$ is $-\tfrac12 c\,\omega^2$, and the pairs of logarithms contribute
$-\log(1 + \omega^2/\theta^2)$ each by the first identity, which is twice the modulus identity;
so the \Levy{} exponent, minus Thorin's, is $\tfrac12 c\,\omega^2 + \tfrac12\int_{(0,\infty)}
\log(1 + \omega^2/\theta^2)\,U(d\theta)$, the right-hand side of \ref{eq:thorin} at
$(\tfrac12 c, U)$.
\end{proof}

The second lemma is the reading of the definition through that computation, and nothing in the
paper uses it. The definition as cited is a moment generating function on real $\sigma$, which the
laws in question do not have at any $\sigma \ne 0$, so what is missing is a formal statement of
the definition at imaginary argument, which is a reading of the source and not a proof.

% shared with blueprint lem:thorin-ggc
\begin{lemma}[the Thorin subclass is the symmetric generalized gamma convolutions]
\label{lem:thorin-ggc}
A symmetric law is an extended generalized gamma convolution in Thorin's sense exactly when its
exponent has the two-sided form of Lemma~\ref{lem:thorin-two-sided} --- that is, when it is
$-\tfrac12c\,\sigma^2 - \int_{\RR\setminus\{0\}}\bigl(\log\tfrac{\theta}{\theta-\sigma}
- \tfrac{\sigma\theta}{1+\theta^2}\bigr)V(d\theta)$ with $c \ge 0$ and $V \ge 0$ symmetric, read
at $\sigma = i\omega$. The two minus signs are the convention used here, in which the exponent is
minus the logarithm of the transform; the source writes the logarithm of the moment generating
function, without them. Consequently, for an admissible $F$ with data $(a,k)$, the two equivalent
conditions of Proposition~\ref{prop:thorin-subclass} hold if and only if the kernels of the
family are symmetric extended generalized gamma convolutions.
\end{lemma}

\begin{proof}
By the cited definition a law belongs to the class exactly when its exponent is
$-b\sigma - \tfrac12 c\sigma^2 - \int\bigl(\log\tfrac{\theta}{\theta-\sigma}
- \tfrac{\sigma\theta}{1+\theta^2}\bigr)V(d\theta)$ for a positive $V$ on $\RR\setminus\{0\}$,
and among these the symmetric ones are those with $b = 0$ and $V$ symmetric. Read at
$\sigma = i\omega$, Lemma~\ref{lem:thorin-two-sided} turns that into
$\tfrac12 c\,\omega^2 + \tfrac12\int_{(0,\infty)}\log(1 + \omega^2/\theta^2)\,U(d\theta)$ with
$U$ the fold of $V$: the compensating terms cancel in pairs, and the pairing of $\pm\theta$
contributes $-\log(1 + \omega^2/\theta^2)$, which is twice the modulus identity there. That is
\ref{eq:thorin} with $a = \tfrac12 c$.

The source's integrability condition on $V$ and \ref{eq:thorin}'s on $U$ are equivalent, and the
comparison is not term by term. The source asks for
$\int_{\RR\setminus\{0\}}(1+\theta^2)^{-1}V(d\theta) < \infty$ and
$\int_{|\theta|\le1}|\log\theta^2|\,V(d\theta) < \infty$; folded, these read
$\int_{(0,\infty)}(1+\theta^2)^{-1}U(d\theta) < \infty$ and
$\int_{(0,1]}2\log(1/\theta)\,U(d\theta) < \infty$. Each direction is an elementary comparison of
integrands. Forwards: on $(0,1]$ one has $\log(1+\theta^{-2}) \le \log 2 + 2\log(1/\theta)$ and
$(1+\theta^2)^{-1} \ge \tfrac12$, so the first condition bounds $U((0,1])$ and the second the
logarithm; on $(1,\infty)$, $\log(1+\theta^{-2}) \le \theta^{-2} \le 2(1+\theta^2)^{-1}$.
Backwards: $\log(1+\theta^{-2}) \ge \log 2$ on $(0,1]$ and
$\ge \tfrac12\theta^{-2} \ge \tfrac12(1+\theta^2)^{-1}$ on $(1,\infty)$, while
$2\log(1/\theta) \le \log(1+\theta^{-2})$ on $(0,1]$. So membership in the class is
Proposition~\ref{prop:thorin-subclass}(2), which is (1) by that proposition.
\end{proof}

\subsection{The Student-t family}\label{sec:student}

Student's t laws are the classical heavy-tailed relatives of the normal law: a normal variable
whose variance is itself random, drawn from an inverse-gamma law. Read through \S\ref{sec:bridge}
that representation is a subordination, Brownian motion run for an inverse-gamma time, and the
inverse-gamma laws are the kernels of the causal Bessel family, so the Student-t family is the
image of that causal corner under the bridge (Proposition~\ref{prop:bridge-families}(3)).
That the t laws are infinitely divisible for every number of degrees of freedom was proved by
Grosswald (\sscite{grosswald1976student}; see \sscite{ismail1979special} for the role of
quotients of Bessel functions). Halgreen proved
that they are self-decomposable, through the corresponding property of the mixing law
(\sscite{halgreen1979self}; \sscite{barndorffnielsen1977infinite} for the mixtures in general),
which is the fact that makes them members of the admissible cone. In the scale-space
literature the family appears at one point. With one degree of freedom the kernel is the Cauchy
kernel, which is the Poisson scale space of Felsberg and Sommer (\sscite{felsberg2001scale};
\sscite{felsberg2004monogenic}) and the member of the $\alpha$-scale spaces that Duits et
al.\ single out (\sscite{duits2004axioms}); it is also the symmetric stable member of index $1$
of \S\ref{sec:stable}, the point where two corners meet.

The proposition gives the kernel and its transform in closed form, places the family in the
Thorin subclass, and counts its moments. The tail is a power, of order $|x|^{-2a-1}$, and the
shape parameter $a$ is a dial on the number of finite moments, those of order below $2a$,
twice the count of the causal ancestor. So this is the corner of the cone with polynomial tails of any
prescribed order, where the \Matern{} corner has exponential tails and every moment, and the
stable corner fewer than two moments.

% shared with blueprint prop:student-t
\begin{proposition}[the Student-t family]\label{prop:student-t}
For $a > 0$ let $T_1$ be the scaled inverse-gamma law of the causal Bessel family, with density
$p_{T_1}(u) = u^{-a-1}e^{-1/(2u)}/(2^a\Gamma(a))$, and let $X_1 := \BM_{T_1}$. Then:
\begin{enumerate}
\item the kernel at canonical scale $1$ is
$\phi_1(x) = \frac{\Gamma(a+\frac12)}{\sqrt\pi\,\Gamma(a)}\,(1+x^2)^{-a-\frac12}$, the
Student-t law with $2a$ degrees of freedom in the scale in which the density is proportional to
$(1+x^2)^{-a-1/2}$, and $\phi_t = t^{-1}\phi_1(\cdot/t)$;
\item the transform is
$e^{-F(\omega)} = \frac{2^{1-a}}{\Gamma(a)}\,|\omega|^aK_a(|\omega|)$;
\item if $T_1$ is a generalized gamma convolution --- Remark~\ref{rem:student-ggc} --- then the
law is self-decomposable and lies in the symmetric Thorin subclass;
\item $\mathbb{E}|X_t|^n < \infty$ if and only if $n < 2a$, and the tail is
$\phi_t(x) \sim c\,t^{2a}|x|^{-2a-1}$, polynomial.
\end{enumerate}
Throughout, $a$ is the shape parameter of the causal Bessel family --- the letter that family
carries in the literature --- and not the Gaussian coefficient of \ref{eq:sd-profile}, which is a
different number and vanishes for these members, as the proof of clause (3) computes under that
clause's hypothesis.
\end{proposition}

The scaling convention, so that the degrees of freedom come out as $2a$ and not $a$, is fixed
by the density in clause (1) and the rate in $p_{T_1}$, and the two should be read together
whenever the family is quoted. The shape parameter is therefore a dial on the number of finite
moments, and the count is twice the causal one, $n < 2a$ against $n < a$ for the delay. Clause (3) is conditional: that the inverse-gamma law is a
generalized gamma convolution is Halgreen's theorem, recorded as
Remark~\ref{rem:student-ggc} and taken here as a hypothesis, and under it the clause follows
from the bridge on Thorin subclasses. Clause (2) is obtained from clause (1) and the Laplace
transform of the delay law, which after an exponential change of variables is the integral by
which $K_a$ is defined, so that no asymptotic of the Bessel function enters.

\begin{proof}[Proof of the parts that are this paper's own]
(1) Conditioning on $T_1$, $\phi_1(x) = \int_0^\infty g_u(x)\,p_{T_1}(u)\,du$; the substitution
$w = 1/u$ turns this into
\[
  \frac{1}{\sqrt{2\pi}\,2^a\Gamma(a)}\int_0^\infty w^{a-\frac12}e^{-w(1+x^2)/2}\,dw
  = \frac{\Gamma(a+\frac12)}{\sqrt{2\pi}\,2^a\Gamma(a)}\Bigl(\frac{2}{1+x^2}\Bigr)^{a+\frac12}
  = \frac{\Gamma(a+\frac12)}{\sqrt\pi\,\Gamma(a)}\,(1+x^2)^{-a-\frac12}.
\]
The scaling clause is Proposition~\ref{prop:bridge-families}. (2) is the causal transform
$e^{-F_{\mathrm I}(\sigma)} = \frac{2^{1-a}}{\Gamma(a)}(\sqrt{2\sigma})^aK_a(\sqrt{2\sigma})$
evaluated at $\sigma = \omega^2/2$; the checked proof computes that Laplace integral rather than
quoting it, and the computation is one substitution: with $A = 1/2$, $B = \sigma$ and
$u = c\,e^{v}$, $c = \sqrt{A/B}$, the integral
$\int_0^\infty u^{-a-1}e^{-1/(2u) - \sigma u}\,du$ becomes
$2c^{-a}\int_0^\infty\cosh(av)\,e^{-2\sqrt{AB}\cosh v}\,dv$, which is the integral by which
$K_a$ is defined here. No asymptotic of $K_a$ enters, at this or any order.

\emph{(3).} Let $U_{\mathrm I}$ be a Thorin measure of $T_1$, as the hypothesis provides:
$\mathbb{E}e^{-\sigma T_1}
= \exp\{-\int_{(0,\infty)}\log(1 + \sigma/\theta)\,U_{\mathrm I}(d\theta)\}$ for $\sigma \ge 0$,
with $\int_{(0,\infty)}\log(1 + \theta^{-1})\,U_{\mathrm I}(d\theta) < \infty$. That exponent is
causally admissible with drift $0$ and delay profile
$k_{\mathrm I}(u) = \int e^{-\theta u}\,U_{\mathrm I}(d\theta)$: a Laplace transform is
nonnegative and nonincreasing, and the causal Frullani identity
$\int_0^\infty(1 - e^{-\sigma u})e^{-\theta u}\,du/u = \log(1 + \sigma/\theta)$ with Tonelli
turns the exponent of Definition~\ref{def:causal-admissible} into the Thorin integral, so that
the two integrability windows follow from its finiteness at $\sigma = 1$ --- for
$1 \wedge u \le e\,(1 - e^{-u})$ bounds the windows by the exponent there. By
Lemma~\ref{lem:bridge-exponents} the image $F_{\mathrm I}(\omega^2/2)$ is admissible, and by the
mixture identity with clause (1) it is an exponent of the Student-t family: the law is
self-decomposable. Complete monotonicity of the spatial profile --- membership in the symmetric
Thorin subclass by Proposition~\ref{prop:thorin-subclass}(1) iff (2) --- is
Proposition~\ref{prop:thorin-subclass}(3) read at that datum, which also returns the Thorin
measure, twice the image of $U_{\mathrm I}$ under $\theta \mapsto \sqrt{2\theta}$, and the
vanishing Gaussian coefficient.

\emph{(4).} The
tail of the density in (1) is $|x|^{-2a-1}$, so $\mathbb{E}|X|^n < \infty$ exactly when
$n < 2a$, which also follows from $\mathbb{E}|\BM_T|^{2n} = (2n-1)!!\;\mathbb{E}T^n$ and
$\mathbb{E}T_1^n < \infty$ exactly when $a > n$.
\end{proof}

% blueprint rem:student-ggc, copied; the cited fact's page anchors are Halgreen's own
\begin{remark}[the inverse-gamma delay law is a generalized gamma convolution]
\label{rem:student-ggc}
The hypothesis that Lemma~\ref{lem:student-subordinated} and clause (3) above quantify over is
a theorem of the literature, and it is the one thing about this corner that the paper does not
prove. For $a > 0$ the inverse-gamma law $T_1$ of shape $a$ is a generalized gamma convolution:
there is a positive measure $U_{\mathrm I}$ on $(0,\infty)$ with
$\int_{(0,\infty)}\log(1 + \theta^{-1})\,U_{\mathrm I}(d\theta) < \infty$ and
\begin{equation}\label{eq:student-ggc}
  \mathbb{E}e^{-\sigma T_1}
   = \exp\Bigl\{-\int_{(0,\infty)}\log\Bigl(1 + \frac{\sigma}{\theta}\Bigr)
       U_{\mathrm I}(d\theta)\Bigr\}, \qquad \sigma \ge 0 .
\end{equation}
That is \sscite{halgreen1979self}, \S2, pp.~14--15, read at $\lambda = -a < 0$, $\chi = 1$,
$\psi = 0$, the parameters at which the generalized inverse Gaussian density of his (1) is the
density $p_{T_1}$ of Proposition~\ref{prop:student-t}. His (5) exhibits the measure: it has
density $u(y) = g_a(2y)$ on $(0,\infty)$, where
$g_\nu(x) = 2\{\pi^2x(J_\nu^2(\sqrt x) + Y_\nu^2(\sqrt x))\}^{-1}$ is Grosswald's
(\sscite{grosswald1976student}), and his integrability conditions (3) ---
$\int_0^1|\log y|\,U_{\mathrm I}(dy) < \infty$ and $\int_1^\infty y^{-1}U_{\mathrm I}(dy) <
\infty$ --- are the displayed finiteness for a measure finite on the compact subsets of
$(0,\infty)$; the displayed form is what finiteness of the exponent at $\sigma = 1$ says, and it
is the form both consumers read. Being a generalized gamma convolution, $T_1$ is in particular
self-decomposable, which is the weaker hypothesis Lemma~\ref{lem:student-subordinated}
quantifies over.

Nothing in the paper depends on this remark: both consumers quantify over
\ref{eq:student-ggc} or over its consequence. The measure Halgreen exhibits is built from the
Bessel functions $J_\nu$ and $Y_\nu$, and no route to \ref{eq:student-ggc} through the
elementary theory is known to the author.
\end{remark}

Which principle cuts the Student-t family out of the cone is a question about locality, and it
is the subject of the next module. Nothing is stated about it here.

\subsection{The symmetric stable family}\label{sec:stable}

The symmetric stable laws are the laws with exponent $c\,|\omega|^\alpha$, $0 < \alpha \le 2$
(\sscite{sato1999levy}, Theorem~14.14, p.~86, for $\alpha < 2$, the Gaussian being the case
$\alpha = 2$): the laws that reproduce themselves under
addition and rescaling, and the attractors of the generalized central limit theorem. They are
the earliest non-Gaussian scale spaces in the axiomatic literature. Pauwels et al.\ obtained
the kernels with transform
$e^{-t|\omega|^\alpha}$, for every $\alpha > 0$, from recursivity and scale invariance, those
with $\alpha \le 2$ being the positive ones (\sscite{pauwels1995extended}). Duits et
al.\ axiomatized them as the $\alpha$-scale spaces, generated by the fractional powers
$-(-\Delta)^{\alpha/2}$ of the Laplacian in the present index, with the Poisson scale space
$\alpha = 1$ as the case they single out (\sscite{duits2004axioms};
\sscite{felsberg2004monogenic}), and Pedersen et al.\ read them as stable \Levy{} motions and
compared them on natural images (\sscite{pedersen2005alpha}). In the hemigroup theory the
family is the semigroup corner: Corollary~\ref{cor:semigroup-case} says that it is exactly what
the axioms give when the stages are required to depend only on the difference of their scales.

The proposition adds the coordinates of the family in this paper's terms. The profile is a pure
power, and the pure powers are the profiles that a dilation changes only by a constant, which
is the stationarity of the increments seen in the profile. The family lies in the Thorin
subclass, with a power-law Thorin density, and it is the image under the bridge of the causal
stable family of half the index. And the moment count is what the semigroup costs: for a pure
power, $\int_1^\infty x^{n-1-\alpha}dx$ is finite only for $n < \alpha$, so no member but the
Gaussian has a variance, whereas every non-power profile with $\int_1^\infty x\,k(x)\,dx$
finite has one. So the heavy tails of the stable scale spaces come from the semigroup
requirement: the same axioms without it admit members with every moment finite. The extended
semigroup class beyond $\alpha = 2$ is not in the cone: positivity fails there.

% DIVERGENCE (2026-09-19, the author's decision D5 at the external review): the blueprint's
% proof ends with a paragraph saying which steps are machine-checked and which route is the
% checked one; the paper's copy omits it (section 1.1 owns that account).
% shared with blueprint prop:stable-family
\begin{proposition}[the symmetric stable family]\label{prop:stable-family}
For $0 < \alpha < 2$ the exponent $F(\omega) = |\omega|^\alpha$ is admissible with $a = 0$ and
profile $k(x) = C_\alpha^{-1}x^{-\alpha}$, where
$C_\alpha = \int_0^\infty (1-\cos u)\,u^{-1-\alpha}du$; its symmetric Thorin density is
proportional to $\theta^{\alpha-1}d\theta$; the kernels are the symmetric $\alpha$-stable
densities; $\mathbb{E}|X_t|^n < \infty$ if and only if $n < \alpha$; and the family is the only
admissible one-parameter semigroup. It is the image under the bridge of the causal stable family
of index $\alpha/2$. The Gaussian ray closes the family at $\alpha = 2$, with $k \equiv 0$ and
$a = 1$.
\end{proposition}

The phrase ``the kernels are the symmetric $\alpha$-stable densities'' names the classical
family whose exponent is $|\omega|^\alpha$. No closed form of the density is asserted, and
outside a few values of $\alpha$ none exists; absolute continuity is
Proposition~\ref{prop:kernel-regularity}. The moment clause is the cited moment criterion of
Proposition~\ref{prop:moments-tails}(1) applied to the pure power.

\begin{proof}
The exponent, the profile and the constant are Corollary~\ref{cor:semigroup-case}, as is the
statement that this is the only admissible one-parameter semigroup. The moment criterion applied
to $k(x) = C_\alpha^{-1}x^{-\alpha}$ gives $\int_1^\infty x^{n-1-\alpha}dx < \infty$ exactly
when $n < \alpha$ (Proposition~\ref{prop:moments-tails}(1)). The causal stable family of index
$\alpha/2$ has $F_{\mathrm I}(\sigma) = \sigma^{\alpha/2}$, whose image is
$2^{-\alpha/2}|\omega|^{\alpha}$ by Proposition~\ref{prop:bridge-families}(4), a positive
multiple of $F$; the causal Thorin density is proportional to $\theta^{\alpha/2-1}d\theta$, and
its image under $\theta' = \sqrt{2\theta}$, with the Jacobian, is proportional to
$\theta'^{\alpha-1}d\theta'$ by Proposition~\ref{prop:thorin-subclass}(3).

The Thorin density is also reached directly, without the bridge and without Bernstein's
theorem: for $c > 0$ the Gamma integral
$\int_0^\infty e^{-\theta x}c\theta^{\alpha-1}d\theta = c\Gamma(\alpha)x^{-\alpha}$ exhibits the
stable profile as a Laplace transform at $c = (C_\alpha\Gamma(\alpha))^{-1}$, and the Tonelli of
Proposition~\ref{prop:thorin-subclass} turns that into \ref{eq:thorin}. Only the positivity of
$C_\alpha$ is used, not its value: the integrand is nonnegative, integrable --- at most
$u^{1-\alpha}/2$ near the origin and at most $2u^{-1-\alpha}$ beyond $1$ --- and nonzero on
$(0,1)$.
\end{proof}
